\section{Deux autres axes de quantification : le moment et la profondeur}
\label{sec:quant-axes}

Les sections précédentes font varier la \emph{calibration} du kernel à format fixe. Deux autres degrés de
liberté commandent le résultat : \emph{quand} on quantifie et \emph{combien de bits} on garde. Les trois
placements comparés ici et les paramètres d'un schéma sub-INT8 sont définis en
\S\ref{sec:method}.

\subsection{Moment : avant, après, ou les deux}

Trois placements sont possibles pour la même tête EWC. \emph{Avant} l'entraînement (QAT, fake-quant dans la
boucle) ; \emph{après} (PTQ sur un modèle FP32 déjà entraîné) ; et \emph{les deux}, c'est-à-dire un
entraînement QAT dont les poids sont ensuite exportés vers le kernel entier calibré --- seul ce dernier
chemin est fidèle au déploiement réel.

En émulation, les trois moments préservent l'\gls{AUROC} sur Monitoring (\num{0.974477} avant, \num{0.974649}
après, \num{0.974663} les deux, contre \num{0.97464} en FP32) : à cette échelle de modèle, aucun moment
n'est disqualifiant. La question tranchée l'est \emph{sur carte}. Le chemin « les deux » a été flashé et
streamé : il donne \num{0.9213} de F1 sur Monitoring et \num{0.9072} sur Pronostia, contre respectivement
\num{0.9173} et \num{0.8995} pour la PTQ calibrée seule (Section~\ref{sec:recovery-board}).

L'écart, \num{0.004} et \num{0.0077}, est réel mais faible. La lecture honnête est donc que
\textbf{la calibration du noyau récupère l'essentiel, et que le QAT n'ajoute pas de gain décisif} sur cette
tête : quand le pipeline d'entraînement ne peut pas être modifié, une PTQ correctement calibrée suffit. Ce
constat contredit la lecture répandue selon laquelle un effondrement en PTQ appellerait nécessairement un
réentraînement.

\begin{figure}[t]
\centering
\includegraphics[width=0.9\linewidth]{fig6_quant_moment.png}
\caption{Moment de quantification. À gauche, les trois moments \emph{émulés PC} (hachuré) ; à droite,
l'A/B \emph{mesuré sur carte} entre la PTQ calibrée seule et le chemin QAT$\to$PTQ.}
\label{fig:moment}
\end{figure}

\subsection{Profondeur : jusqu'où descendre sous 8 bits}

Sous les 8 bits, la question devient : à partir de quelle profondeur la métrique décroche, et la
granularité du facteur d'échelle change-t-elle cette limite ? Le balayage émulé répond aux deux.

La limite dépend du jeu ; la grille complète est en Annexe~\ref{app:depth}
(Tableau~\ref{tab:annex-depth}). Sur Monitoring, la dégradation reste contenue jusqu'au binaire
($\Delta$AUROC $=-\,\num{0.0117}$). Sur Pronostia, le binaire décroche ($-\,\num{0.0275}$) mais le ternaire
tient ($-\,\num{0.0153}$) : lue en AUROC, la frontière exploitable serait donc le \emph{ternaire}, le
binaire n'étant défendable que sur le jeu le plus facile.

\paragraph{La frontière dépend de la métrique de sélection.} Cette lecture ne survit pas à la confrontation
avec le F1 \emph{mesuré sur carte} des mêmes schémas, que le Tableau~\ref{tab:depth-board} place à côté de
l'AUROC. Sur Pronostia, le ternaire conserve une AUROC de \num{0.99907} mais son F1 tombe à \num{0.379} ;
le binaire garde une AUROC de \num{0.96616} pour un F1 \emph{nul}, c'est-à-dire un modèle qui ne prédit
plus jamais la classe de défaut. L'ordre des scores reste bon là où le seuil de décision, lui, ne l'est
plus : une métrique de rang est insensible au déséquilibre des classes, un F1 ne l'est pas. Sur les deux
jeux, le seul schéma sub-INT8 qui préserve réellement la métrique d'exploitation est l'\textbf{INT4}
(\num{0.921} et \num{0.9887}).

Nous rapportons les deux lectures plutôt que d'en choisir une, car leur écart est lui-même le résultat :
\textbf{ce n'est pas la profondeur seule qui fixe la frontière, c'est le couple profondeur--métrique de
sélection}. Un travail qui ne rapporterait que l'AUROC conclurait au ternaire et déploierait, sur
Pronostia, un détecteur qui ne détecte rien.

% Généré par scripts/generate_article_tables.py — NE PAS ÉDITER À LA MAIN.
% Source : experiments/exp_S40_article_metrics/summary.json
\begin{table}[htbp]
\centering
\caption{Schémas sub-INT8 \emph{mesurés sur carte} : AUROC et F1 de la classe \texttt{faulty} sur les mêmes cellules. L'AUROC reste haute là où le F1 décroche, la profondeur ne déplaçant pas les deux métriques ensemble.}
\label{tab:depth-board}
\begin{tabular}{lcccc}
\toprule
Profondeur & \multicolumn{2}{c}{Monitoring} & \multicolumn{2}{c}{Pronostia} \\
\cmidrule(lr){2-3}\cmidrule(lr){4-5}
 & AUROC & F1 & AUROC & F1 \\
\midrule
\texttt{int4} & \num{0.98787} & \num{0.9210} & \num{0.99995} & \num{0.9887} \\
\texttt{ternaire} & \num{0.98409} & \num{0.9042} & \num{0.99907} & \num{0.3790} \\
\texttt{binaire} & \num{0.96784} & \num{0.6878} & \num{0.96616} & \num{0.0000} \\
\bottomrule
\end{tabular}
\end{table}


La granularité, elle, repousse le décrochage. À 2 bits sur Pronostia, passer d'une échelle par-tenseur à
une échelle par-canal fait passer $\Delta$AUROC de $-\,\num{0.0464}$ à $-\,\num{0.0086}$. Un point négatif
mérite d'être rapporté aussi : le \emph{zéro-point} affine ne rachète rien --- à 2 bits sur Monitoring il
dégrade même nettement ($-\,\num{0.0619}$ contre $-\,\num{0.0069}$ en symétrique), et le
Tableau~\ref{tab:annex-symmetry} montre que le constat vaut aux trois profondeurs testées.

\paragraph{Le gain RAM n'existe que packé.} Le point crucial est que la RAM annoncée par ces schémas est
\emph{théorique} tant que les poids restent stockés dans un conteneur \texttt{int8}. Mesuré sur carte, le
\texttt{.bss} d'un modèle 4, 2 ou 1 bit \emph{non packé} est strictement identique à celui de l'INT8 : le
gain n'apparaît qu'avec un véritable bit-packing, où il croît quand les bits baissent
(\num{336}, \num{504} puis \num{572}~octets économisés sur Monitoring). Le dépacking se paie en latence,
d'environ \SI{55}{\micro\second} (de \SI{67}{\micro\second} non packé à \SI{123}{\micro\second} packé sur
Monitoring) --- sans menacer le budget Gap~2. La parité contre l'émulateur reste \num{1.000} sur les douze
cellules mesurées : le schéma est porté sans perte.

\begin{figure}[t]
\centering
\includegraphics[width=0.9\linewidth]{fig7_quant_depth_packing.png}
\caption{Profondeur sub-INT8. À gauche, $\Delta$AUROC par profondeur et granularité (\emph{émulé PC}) ---
la lecture en F1 mesuré sur carte, qui déplace la frontière, est au Tableau~\ref{tab:depth-board} ;
à droite, les octets de \texttt{.bss} réellement économisés par le bit-packing (\emph{mesuré sur carte}).}
\label{fig:depth}
\end{figure}
